The Convex Hull Ratio (CH $= A / A(\text{conv}(D))$) is the most visually
intuitive compactness metric for redistricting litigation: it measures what
fraction of a district's convex envelope is actually occupied by the district,
detecting non-convex ``tentacle'' or ``octopus arm'' geometries that perimeter-based
metrics fail to identify.
We study CH across four \textsc{bisect} structure algorithms---standard-bisect,
prime-factor, ratio-optimal, and moving-knife---for North Carolina ($k=14$),
Wisconsin ($k=8$), and Texas ($k=38$) under 2020 Census data.
Key findings: (1) Synthetic tentacle districts with PP $\approx 0.15$
and Reock $\approx 0.35$ score CH $< 0.60$, demonstrating CH's unique
sensitivity to non-convex appendages.
(2) All four bisect structure algorithms produce CH $> 0.85$ on all NC 2020
congressional districts, confirming that bisect's recursive partitioning
does not generate tentacle geometries by construction.
(3) CH is the most legally intuitive compactness metric for courtroom
presentation under \textit{Shaw v.\ Reno} (1993)'s visual-irregularity test.
All results use single seed $= 42^\dagger$.
